\section*{Interleaved disjoint activation is monotone}

All logarithms are natural.  Put
\[
 \psi(t)=t^2\log(1/t^2),\qquad
 K(x,y)=\psi(x)+\psi(y)-\frac{\psi(x+y)+\psi(x-y)}2.
\]

\subsection*{Scalar kernel lemma}

If $|x|+|y|\le1$, then
\[
 0\le K(x,y)\le\psi(y)+3y^2.                         \tag{1}
\]
Both constants are sharp in their stated form.

\paragraph{Proof of the lower bound.}
The kernel is unchanged by changing either sign and is homogeneous of
degree two.  It is therefore enough to take $x,y\ge0$ and $x+y=1$.
Write $t=x-y\in[0,1]$, so $x=(1+t)/2$, $y=(1-t)/2$, and define
\[
 L(t)=2\psi((1+t)/2)+2\psi((1-t)/2)-\psi(t)=2K(x,y).
\]
Since $\psi'(r)=-4r\log r-2r$ for $r>0$,
\[
 L'(t)=4\{t\log t-x\log x+y\log y\}\le0.
\]
Indeed $x=t+y$ and
$x\log x-t\log t-y\log y=xH(t/x)\ge0$.  As $L(1)=0$, the lower bound
follows.

For the upper bound, the function $G(r)=\psi(r)+3r^2$ is convex on
$[-1,1]$, because $G''(r)=2\log(1/r^2)\ge0$.  Hence
\[
 \frac{\psi(x+y)+\psi(x-y)}2+3(x^2+y^2)
 \ge\psi(x)+3x^2,
\]
which rearranges to the upper bound in (1).  The coefficient $3$ is
approached as $x\uparrow1$ and $y\downarrow0$. \(\square\)

\subsection*{The interleaved activation theorem}

Let $W$ and $X$ be disjoint physical coordinate blocks, let $a,b$ be
functions on $W$, let $h$ be Boolean on $X$, and assume $a+b$ and $a-b$
are Boolean.  Define
\[
 F=a+hb,\qquad q=\|b\|_2^2,
 \qquad \mathcal J(a,b;h)=\mathcal E_{\rm raw}(F)
                         -q\mathcal E_{\rm raw}(h).                \tag{2}
\]

Consider any adaptive restriction tree which may interleave queries from
$W$ and $X$, never querying a physical coordinate twice on one path, and
which is complete on the source block $X$.  At a source-query node $v$,
let $R_v$ be the exact biased coherent remainder for the current restricted
pair $(a_v,b_v)$ and source $h_v$.  Then
\[
 \boxed{\sum_{v:\,i(v)\in X}w_vR_v\le\mathcal J(a,b;h).}          \tag{3}
\]
More precisely,
\begin{align*}
 \sum_{v:\,i(v)\in X}w_vR_v
  ={}&\mathcal J(a,b;h)
     -\sum_{\ell\ {
m leaf}}w_\ell\mathcal J(a_\ell,b_\ell;h_\ell)\\
    &-\sum_{v:\,i(v)\in W}w_vD_{i(v)}(F_v),                       \tag{4}
\end{align*}
and both subtracted terms are nonnegative.

\paragraph{Proof.}
At a source node, the fixed-pair identity gives
\[
 R_v=\mathcal J(a_v,b_v;h_v)
   -\mathbb E[\mathcal J(a_v,b_v;h_{\rm child})\mid v].           \tag{5}
\]
At an activation node $z\in W$, the source is unchanged and Parseval gives
$q_v=(q_{v+}+q_{v-})/2$.  Therefore the Doob decrement of (2) is exactly
\[
 \mathcal J(a_v,b_v;h_v)
 -\mathbb E[\mathcal J(a_{\rm child},b_{\rm child};h_v)\mid v]
 =D_z(F_v).                                                        \tag{6}
\]
If $F_v=P+zQ$, then
\[
 D_z(F_v)=\sum_TK(\widehat P(T),\widehat Q(T)).                    \tag{7}
\]
The two child coefficients are $\widehat P(T)\pm\widehat Q(T)$ and have
absolute value at most one, so
$|\widehat P(T)|+|\widehat Q(T)|\le1$.  The kernel lemma shows that every
term in (7) is nonnegative.  Summing the Doob decrements over the whole
tree yields (4).  At a leaf $h_\ell=\pm1$, hence
\[
 \mathcal J(a_\ell,b_\ell;h_\ell)
 =\mathcal E_{\rm raw}(a_\ell\pm b_\ell)\ge0.
\]
Dropping the two nonnegative terms proves (3). \(\square\)

The upper half of (1) also gives the exact activation-channel envelope
\[
 D_z(F_v)\le\mathcal E_{\rm raw}(Q)+3\|Q\|_2^2,                   \tag{8}
\]
where $Q=(F_{v,+}-F_{v,-})/2$.  Thus the activation loss has a derivative
entropy plus influence form, although only its nonnegativity is needed in
(3).

For one activation bit, write $a=A+zC$ and $b=B+zD$.  If
$\mu=\mathbb Eh$, Fourier separation of the disjoint blocks and homogeneity
of $K$ give the more explicit formula
\[
 D_z(F)=\sum_S K(A_S+\mu B_S,C_S+\mu D_S)
       +(1-\mu^2)\sum_S K(B_S,D_S).                    \tag{9}
\]
Thus the entire two-state activation mismatch depends on the source only
through its mean and variance, not through its dimension or spectral shape.

\subsection*{Selector corollary}

For the one-bit selector pair
\[
 a(x)=\frac{1+x}{2},\qquad b(x)=\frac{1-x}{2},qquad h=2s-1,
 \qquad p=\mathbb Es,
\]
the root defect is
\[
 Q(p)=\psi(p)+\psi(1-p)-\frac12\psi(2p-1)+4p(1-p)\log2.
\]
Consequently every adaptive tree which arbitrarily interleaves the selector
query with all source queries satisfies
\[
 \sum_{v:\,i(v)\in X}w_vR_v
 \le Q(p)\le8p\log2\le32\log2\,I\!\left(\frac{1-x}{2}s\right).   \tag{10}
\]
This is uniform in dimension, depth, query order, bias, and coherent
collision ratios.  Querying the selector first spends the root defect in
the nonnegative activation decrement; querying the source first can retain
the full endpoint $Q(p)$.

\subsection*{Scope}

The theorem handles genuinely changing activation weights caused by
restrictions of a disjoint activation block.  It does not cover a physical
coordinate that simultaneously changes the activation pair and the source,
nor reused-selector overlap.  In that shared-coordinate regime the
Parseval cancellation in (6) acquires the crossed terms from the four-channel
identity, and their multi-query payment remains open.

